\scanpage{15}
coh. de De Rham de $X/S$
\[
H^*_{\mathrm{DR}}(X/S,\underline V)\simeq H^*_{\mathrm{str}}(X/S,\underline V)
\]
au second mb duquel il n'est plus question de \deleted{cohomologie} formes différentielles !

Si on ne fait plus d'hyp. de lissité ni de caractéristique, il reste vrai qu'on a un hom canonique \inserted{compatible aux stratifications}
\diagram{\begin{tikzcd}[ampersand replacement=\&]
\underline V\arrow[r]\&\mathcal H^0(\ell Q^0[\Omega^\bullet(V)])
\end{tikzcd}}
i.e. un hom de complexes de Modules stratifiés
\diagram{\begin{tikzcd}[ampersand replacement=\&]
\underline V\arrow[r]\&\deleted{\unsure{\ell}}Q^0[\Omega^\bullet(V)]
\end{tikzcd}}
d'où en passant aux $C^\bullet_{\mathrm{str}}$ des deux \unsure{mbs}, un \uline{hom. can.}
\diagram{\[\boxed{\begin{tikzcd}[ampersand replacement=\&,column sep=large]
H^*(X_{\inserted{\mathrm{str}}}\deleted{\unsure{/S}},\underline V)\arrow[r]\&H^*_{\mathrm{DR}}(X/S,\underline V)
\end{tikzcd}}\]}
\margin{14}
